\chapter{Mínimos Quadrados Ordinários}
\label{cap:ols}

% ─────────────────────────────────────────────────────────────────────────────
\section{O modelo}

Considere o modelo de regressão linear:

\[
  y = X\beta + \varepsilon
\]

onde $y$ é um vetor $n \times 1$ de observações da variável dependente,
$X$ é uma matriz $n \times k$ de regressores (a primeira coluna é
identicamente 1, correspondendo ao intercepto), $\beta$ é um vetor
$k \times 1$ de parâmetros desconhecidos, e $\varepsilon$ é um vetor
$n \times 1$ de erros.

\subsection*{Hipóteses de Gauss-Markov}

\begin{enumerate}
  \item \textbf{Linearidade:} $y = X\beta + \varepsilon$.
  \item \textbf{Posto completo:} $\mathrm{rank}(X) = k$ (sem multicolinearidade perfeita).
  \item \textbf{Exogeneidade estrita:} $E[\varepsilon \mid X] = 0$.
  \item \textbf{Homocedasticidade:} $\mathrm{Var}[\varepsilon \mid X] = \sigma^2 I_n$.
  \item \textbf{Não-autocorrelação:} $\mathrm{Cov}[\varepsilon_i, \varepsilon_j \mid X] = 0$ para $i \neq j$.
\end{enumerate}

Sob H1–H3 o OLS é não-viesado. Sob H1–H5 é BLUE (Teorema de Gauss-Markov).

% ─────────────────────────────────────────────────────────────────────────────
\section{O estimador}

O estimador OLS minimiza a soma dos quadrados dos resíduos:

\[
  \hat{\beta} = \arg\min_{\beta}\, (y - X\beta)'(y - X\beta)
\]

A condição de primeira ordem $X'(y - X\beta) = 0$ — as \emph{equações normais} —
tem solução analítica:

\[
  \boxed{\hat{\beta} = (X'X)^{-1}X'y}
\]

Os resíduos são $\hat{e} = y - X\hat{\beta} = M_X y$, onde
$M_X = I - X(X'X)^{-1}X'$ é a matriz de resíduos (simétrica e idempotente).

% ─────────────────────────────────────────────────────────────────────────────
\section{Propriedades}

Sob H1–H3:
\[
  E[\hat{\beta}] = \beta
\]

Sob H1–H5, a variância condicional do estimador é:
\[
  \mathrm{Var}[\hat{\beta} \mid X] = \sigma^2 (X'X)^{-1}
\]

O estimador não-viesado da variância do erro é:
\[
  s^2 = \frac{\hat{e}'\hat{e}}{n - k}
\]

e portanto $\widehat{\mathrm{Var}}[\hat{\beta}] = s^2 (X'X)^{-1}$.

% ─────────────────────────────────────────────────────────────────────────────
\section{Erros padrão}

\subsection*{Clássico}

Assume homocedasticidade. A matriz de covariância estimada é
$s^2(X'X)^{-1}$ e o erro padrão do coeficiente $j$ é
$\mathrm{se}_j = s\,\sqrt{[(X'X)^{-1}]_{jj}}$.

\subsection*{Robusto à heterocedasticidade}

Quando $\mathrm{Var}[\varepsilon_i \mid x_i] = \sigma_i^2$ (heterocedasticidade),
o estimador sanduíche fornece inferência válida.

\textbf{HC1} (White com correção de graus de liberdade):
\[
  \hat{V}_{\mathrm{HC1}} = \frac{n}{n-k}\,(X'X)^{-1}
    \!\left(\sum_{i=1}^{n} \hat{e}_i^2\, x_i x_i'\right)\!(X'X)^{-1}
\]

\textbf{HC3} (mais conservador, recomendado em amostras pequenas):
\[
  \hat{V}_{\mathrm{HC3}} = (X'X)^{-1}
    \!\left(\sum_{i=1}^{n} \frac{\hat{e}_i^2}{(1 - h_{ii})^2}\, x_i x_i'\right)\!(X'X)^{-1}
\]

onde $h_{ii} = x_i'(X'X)^{-1}x_i$ é o \emph{leverage} da observação $i$.

\subsection*{Clusterizado}

Quando os erros são correlacionados dentro de grupos $g = 1,\ldots,G$:
\[
  \hat{V}_{\mathrm{cl}} = \frac{G(n-1)}{(G-1)(n-k)}\,
    (X'X)^{-1}\!\left(\sum_{g=1}^{G} X_g'\hat{e}_g\hat{e}_g'X_g\right)\!(X'X)^{-1}
\]

% ─────────────────────────────────────────────────────────────────────────────
\section{Inferência}

\subsection*{Teste $t$ sobre um coeficiente}

Sob H0: $\beta_j = \beta_j^0$:
\[
  t_j = \frac{\hat{\beta}_j - \beta_j^0}{\mathrm{se}(\hat{\beta}_j)}
  \;\xrightarrow{d}\; t_{n-k}
\]

\subsection*{Teste $F$ sobre restrições lineares}

Para $R\beta = r$ com $q$ restrições:
\[
  F = \frac{(R\hat{\beta} - r)'
       \bigl[R\,(X'X)^{-1}R'\bigr]^{-1}
       (R\hat{\beta} - r)}{q\,s^2}
  \;\xrightarrow{d}\; F_{q,\,n-k}
\]

\begin{nota}
  Com erros robustos ou clusterizados, $\hat{V}$ substitui $s^2(X'X)^{-1}$
  no cálculo do teste $F$.
\end{nota}

% ─────────────────────────────────────────────────────────────────────────────
\section{Sintaxe}

\begin{lstlisting}
ols(fórmula, df)
ols(fórmula, df, cov=tipo)
ols(fórmula, df, cov=cluster, cluster=coluna)
ols(fórmula, df, if = condição)
\end{lstlisting}

Alias: \lstinline{reg()} é equivalente a \lstinline{ols()}.

\section{Fórmula}

A fórmula segue a notação padrão \texttt{y \textasciitilde{} x1 + x2 + ...}:

\begin{lstlisting}[caption={OLS básico}]
input df
Y X
10 2
12 3
8  1
15 5
11 2
14 4
end

ols(Y ~ X, df)
\end{lstlisting}

A constante é incluída automaticamente. Para múltiplas variáveis:

\begin{lstlisting}
ols(Y ~ X1 + X2 + X3, df)
\end{lstlisting}

A fórmula pode ser passada como string:

\begin{lstlisting}
let formula = "Y ~ X1 + X2"
ols(formula, df)
\end{lstlisting}

\section{Opções de erro padrão}

\begin{center}
\begin{tabular}{ll}
\toprule
\textbf{Opção} & \textbf{Fórmula} \\
\midrule
(padrão)                            & $s^2(X'X)^{-1}$ \\
\texttt{cov=robust}                 & HC1 (sanduíche White) \\
\texttt{cov=HC1}                    & idem \\
\texttt{cov=HC3}                    & HC3 (leverage-ajustado) \\
\texttt{cov=cluster, cluster=col}   & clusterizado por \texttt{col} \\
\bottomrule
\end{tabular}
\end{center}

\begin{lstlisting}[caption={Variantes de erro padrão}]
ols(Y ~ X, df)                           // clássico
ols(Y ~ X, df, cov=robust)              // HC1
ols(Y ~ X, df, cov=HC3)                 // HC3
ols(Y ~ X, df, cov=cluster, cluster=id) // clusterizado
\end{lstlisting}

\section{Amostra condicional}

A opção \lstinline{if =} restringe a estimação a um subconjunto sem modificar o DataFrame:

\begin{lstlisting}
ols(Y ~ X, df, if = grupo == 1)
\end{lstlisting}

\section{Capturando o resultado}

\begin{lstlisting}
let m = ols(Y ~ X, df)
display m
\end{lstlisting}

Quando atribuído, \lstinline{ols()} não imprime a tabela — apenas armazena o modelo.

\section{Pós-estimação}

\subsection{\texttt{predict}}

\begin{lstlisting}
let m = ols(Y ~ X, df)

let yhat = predict(m, "xb")        // $\hat{y} = X\hat{\beta}$
let ehat = predict(m, "residuals") // $\hat{e} = y - X\hat{\beta}$
\end{lstlisting}

Aliases aceitos: \texttt{"fitted"} para valores ajustados; \texttt{"resid"} ou \texttt{"e"} para resíduos.

\subsection{\texttt{vif}}

O VIF do regressor $j$ é $\mathrm{VIF}_j = 1/(1 - R_j^2)$, onde $R_j^2$ é o
$R^2$ da regressão de $x_j$ sobre os demais regressores. Valores acima de 10
indicam multicolinearidade severa.

\begin{lstlisting}
let m = ols(Y ~ X1 + X2 + X3, df)
vif(m)
\end{lstlisting}

\subsection{\texttt{test}}

Teste $F$ de restrições lineares $H_0: R\beta = r$:

\begin{lstlisting}
let m = ols(Y ~ X1 + X2 + X3, df)
test(m, X2 = 0, X3 = 0)    // H0: X2 = X3 = 0
\end{lstlisting}

\subsection{\texttt{lincom}}

Combinação linear $c'\hat{\beta}$ com erro padrão $\sqrt{c'\hat{V}c}$:

\begin{lstlisting}
let m = ols(Y ~ X1 + X2, df)
lincom(m, X1 + X2)
\end{lstlisting}

\subsection{\texttt{margins}}

Efeitos marginais médios. No OLS linear $\partial E[y]/\partial x_j = \beta_j$,
equivalente aos coeficientes.

\begin{lstlisting}
let m = ols(Y ~ X, df)
margins(m)
\end{lstlisting}

\section{Tabela de resultados com \texttt{esttab}}

\begin{lstlisting}[caption={Comparação de especificações}]
let m1 = ols(Y ~ X, df)
let m2 = ols(Y ~ X, df, cov=robust)
let m3 = ols(Y ~ X1 + X2, df)

esttab(m1, m2, m3)
\end{lstlisting}

\begin{lstlisting}
eststo(ols(Y ~ X, df))
eststo(ols(Y ~ X1 + X2, df))
esttab()
\end{lstlisting}

\section{Bootstrap}

Erros padrão por bootstrap — alternativa não-paramétrica:

\begin{lstlisting}
bootstrap(ols, Y ~ X, df, n=1000)
bootse(m, n=500)
\end{lstlisting}
